\subsection{Анализ подходов к программированию}

При разработке скриптов для дрона можно использовать разные архитектурные паттерны. Рассмотрим три основных подхода.

\subsubsection{1. Блокирующее выполнение (Sleep)}
Самый простой, но \textbf{не рекомендуемый} способ. Программа выполняет действие, затем «засыпает» на определенное время с помощью функции \texttt{sleep()}, затем выполняет следующее действие.

\begin{warnbox}[Почему sleep — это плохо?]
На борту квадрокоптера функция \texttt{sleep()} останавливает выполнение всего потока скрипта. Это может помешать своевременной реакции на датчики, команды управления или аварийные ситуации. Используйте этот подход только для простейших тестов на земле.
\end{warnbox}

\subsubsection{2. Асинхронные таймеры}
Это «золотой стандарт» программирования встраиваемых систем. Вместо того чтобы ждать, мы поручаем системе вызвать нашу функцию через определённый интервал времени.

\begin{infobox}[Преимущества таймеров]
\begin{itemize}
    \item \textbf{Параллельность:} можно запустить несколько независимых анимаций (например, мигание навигационных огней и перелив основной ленты).
    \item \textbf{Отзывчивость:} скрипт не блокируется и готов обрабатывать другие события.
\end{itemize}
\end{infobox}

Основной инструмент — \texttt{Timer.new(period, fn)}. Функция \texttt{fn} должна быть короткой и выполнять только один шаг анимации (тик).

\subsubsection{3. Конечный автомат (State Machine)}
Подход для сложной логики. Программа всегда находится в одном из состояний (например, \texttt{IDLE}, \texttt{TAKEOFF}, \texttt{ALARM}). Переходы между состояниями происходят по событиям.

Этот подход идеально подходит для миссий, где поведение дрона и светодиодов зависит от внешних условий (взлетел, сел, разрядился).
